% !TeX root = ../nc-cubic-unmark.tex
\section{構造対象と $\mathcal O$-marking}\label{sec:weak-definition}
本節では $\kk=\mathbb C$ とする。
非可換 del Pezzo の内在的定義に現れる構造対象の選択を
$\mathcal O$-marking と呼ぶ。
この定義と、後の線束・Weyl 群の定理で用いる追加仮定を分ける。

\begin{definition}[$\mathcal O$-marked del Pezzo と cubic]
\label{def:o-marked-cubic}
$d\in\{1,\ldots,9\}$ とする。
$\mathcal A$ は $\kk$ 線型 noetherian Abel 圏で、$D^b(\mathcal A)$ は
Hom 有限で Serre 関手 $\mathsf S$ を持ち、
$H=\mathsf S^{-1}[2]$ が $H(\mathcal A)=\mathcal A$ を満たすとする。
$\mathcal A_0$ を有限長対象の Serre 部分圏、
$\mathcal A_1$ を $\mathcal A/\mathcal A_0$ で有限長になる対象の逆像とし、
$q:\mathcal A\to\mathcal A/\mathcal A_1$ を商関手とする。
対象 $\OO\in\mathcal A$ が\emph{次数 $d$ の構造対象}であるとは、次を満たすことをいう。
\begin{enumerate}
\item $q\OO$ は非零単純対象であり、
      $\Hom(T,\OO)=0$ がすべての $T\in\mathcal A_1$ に対して成り立つ。
      これを Gabriel--Krull の意味で rank-one torsion-free と呼ぶ。
\item $(\OO,H)$ は Artin--Zhang 豊富である。
      すなわち全対象は十分負の $H$-twist の有限直和の商となり、
      任意の全射 $M\twoheadrightarrow N$ に対して
      $\Hom(\OO,H^nM)\to\Hom(\OO,H^nN)$ は $n\gg0$ で全射となる。
\item 反標準切断環 $R_{\OO}=\bigoplus_{n\ge0}\Hom(\OO,H^n\OO)$ は
      次の Hilbert 級数を持つ。
      \begin{equation}\label{eq:o-del-pezzo-hilbert}
       h_{R_{\OO}}(t)=\frac{1+(d-2)t+t^2}{(1-t)^3},\qquad
       \dim\Hom(\OO,H^n\OO)=1+\frac d2n(n+1)\quad(n\ge0).
      \end{equation}
\end{enumerate}
組 $(\mathcal A,\OO)$ を\emph{$\mathcal O$-marked nc del Pezzo of degree $d$}、
$d=3$ の場合を\emph{$\mathcal O$-marked nc cubic}と呼ぶ。
基礎圏 $\mathcal A$ がこのような対象を少なくとも一つ持つことを、
本稿で採用する無標識 nc del Pezzo の内在的条件とする。
構造対象の例外性、smooth proper dg enhancement、数値 Picard 格子、
AS 表示および ambient 非可換射影空間は、この定義には含めない。
\end{definition}

ここで ``structure sheaf'' はテンソル単位を意味しない。
対象間のテンソル積を仮定せず、同じ $\mathcal A$ の中に複数の構造対象があってよい。
Serre 関手は圏から定まり、$H$ は別の自由な標識データではない。
exact 同値は Gabriel--Krull filtration と Serre 双対性を輸送するので、
上の条件は $(\mathcal A,\OO)$ の exact 同値で不変である。

\begin{lemma}[最小公理から従う消滅と次元]
\label{lem:o-minimal-consequences}
定義\ref{def:o-marked-cubic}の下で、$\mathcal A$ の homological 次元は $2$ である。
また
\begin{equation}\label{eq:o-minimal-vanishing}
 \End(\OO)=\kk,\qquad
 \Hom(\OO,H^{-b}\OO)=0\ (b>0),\qquad
 \Ext^2(\OO,H^n\OO)=0\ (n\ge0).
\end{equation}
$\mathcal A/\mathcal A_1$ の全対象は有限長であり、その単純対象は
$q\OO$ の一つの同型類だけである。
\end{lemma}
\begin{proof}
Serre 双対性から
$\Ext^i(M,N)\simeq\Hom(HN,M[2-i])^*$ となり、$i>2$ では零である。
非零 $M$ に対して $N=H^{-1}M$ と取ると
$\Ext^2(M,N)\simeq\End(M)^*\ne0$ なので次元はちょうど $2$ である。
$\End(\OO)=\kk$ は Hilbert 関数の $n=0$ からも従う。

rank-one torsion-free 対象間の非零射 $f:E\to F$ は単射である。
実際、$qf=0$ なら $\operatorname{im}f\in\mathcal A_1$ が $F$ に埋め込まれて矛盾し、
$qf\ne0$ なら単純対象間の同型なので $\ker f\in\mathcal A_1$ となって核は零である。
$H$ は exact だから、この議論は $H$-twist にも適用できる。
もし $\OO\to H^{-b}\OO$ が非零なら単射であり、
$h(n)=\dim\Hom(\OO,H^n\OO)$ に対して $h(n)\le h(n-b)$ が全整数 $n$ で成り立つ。
各剰余類 $n\bmod b$ でこれは $h(n)$ の有界性を与え、
式\eqref{eq:o-del-pezzo-hilbert}の二次増大と矛盾する。
第二 Ext の消滅は Serre 双対性とこの負の Hom 消滅による。

$j\ge i$ なら $\Hom(H^i\OO,H^j\OO)\ne0$ であり、
その非零射は商圏の単純対象間の同型を与える。
従って全 $qH^i\OO$ は $q\OO$ に同型である。
豊富性による有限生成全射を $q$ で写せば、任意の $qM$ は
$q\OO$ の有限直和の商となる。よって結論を得る。
\end{proof}

\begin{lemma}[中心と自然自己同型]\label{lem:dimension-centre}
定義\ref{def:o-marked-cubic}の下で
\[
 \operatorname{Nat}(\mathrm{id}_{\mathcal A},\mathrm{id}_{\mathcal A})=\kk.
\]
恒等関手の可逆自然変換は $\kk^\times$ であり、$\OO$ 上の恒等を要求すれば自明である。
同じことは、構造対象上の同型を指定した exact 同値の自然自己同型にも成り立つ。
\end{lemma}
\begin{proof}
自然変換 $\eta$ は各 $H^n\OO$ 上でスカラー $c_n$ として作用する。
$\Hom(H^i\OO,H^j\OO)\ne0$ ($j\ge i$) に対する自然性から $c_i=c_j$ である。
任意の対象は有限個の負 twist の直和の商なので、$\eta$ は全対象上で同じスカラーである。
exact 同値の場合はその準逆を合成し、恒等関手の場合に帰着する。
\end{proof}

\begin{proposition}[構造対象による反標準表示]
\label{prop:o-anticanonical-reconstruction}
$\mathcal O$-marked cubic $(\mathcal A,\OO)$ に対して、
Artin--Zhang 再構成は
\[
 (\mathcal A,\OO,H)\simeq(\qgr R_{\OO},\pi R_{\OO},(1))
\]
という exact 同値を与える。
従ってここで採用する構造対象は、反標準表示において実際に
$\pi R_{\OO}$ という構造層として現れる。
\end{proposition}
\begin{proof}
$\OO$ は noetherian、$\End(\OO)=\kk$ は右 noetherian、
全 $\Hom(\OO,M)$ は有限次元であり、$(\OO,H)$ は豊富である。
これらは Artin--Zhang\cite[Theorem 4.5]{AZ}の再構成の仮定である。
\end{proof}

\subsection{線束・Weyl 群の比較で用いる強化条件}
\begin{definition}[例外的構造対象と比較用の弱公理]\label{def:weak-cubic}
定義\ref{def:o-marked-cubic}の $d=3$ の組 $(\mathcal A,\OO)$ が、さらに
$D^b(\mathcal A)$ の smooth proper dg enhancement と
\begin{equation}\label{eq:exceptional}
 \Ext^i(\OO,\OO)=0\qquad(i>0)
\end{equation}
を持つとき、本稿の\emph{比較用の弱公理}を満たすと呼ぶ。
後の数値分類でいう許容線束は、この例外的な構造対象である。
この範囲では rank 条件と Hilbert 関数を
\begin{gather}
 q\OO\text{ が非零単純},\qquad
 \Hom(T,\OO)=0\quad(T\in\mathcal A_1),\label{eq:gk-rank}\\
 \dim\Hom(\OO,H^m\OO)=
 \begin{cases}1+\frac32m(m+1)&m\ge0,\\0&m<0\end{cases}
 \label{eq:weak-hilbert}
\end{gather}
と記す。負方向の消滅は補題\ref{lem:o-minimal-consequences}の帰結である。
以下で「弱標識」と記す組は、この比較用の弱公理を満たす
$\mathcal O$-marking を意味し、最小定義全体を指さない。
\end{definition}

最小の $\mathcal O$-marking とこの強化を同一視しない。
前者の公理から $\Ext^1(\OO,\OO)=0$ が従うかは、本稿では未解決である。
後の命題\ref{prop:o-marking-defect-comparison}は、既知の三次曲面型 heart の中で
この差を数値的欠損と自己拡大によって正確に記述する。
$Q$-marking の先頭対象は例外対象なので、その忘却像は例外的な
$\mathcal O$-marked 部分に入る。

\begin{theorem}[固定圏における構造対象の選択と同値類]
\label{thm:forgetful-fibres}
固定した $\mathcal A$ の次数 $3$ の構造対象の同型類集合を
$\operatorname{Str}_3(\mathcal A)$ とする。
基礎圏との同定を保持した $\mathcal O$-marking の選択は
$\operatorname{Str}_3(\mathcal A)$、同定を忘れた組の exact 同値類は
\[
 \operatorname{Aut}_{\rm ex}(\mathcal A)\backslash
 \operatorname{Str}_3(\mathcal A)
\]
である。例外的部分に限れば $\operatorname{Str}_3$ を
$\operatorname{Line}_{\rm adm}$ に置き換える。
固定した構造対象上の $Q$-marking の選択も同様に対象の同型類で考えるが、
その moduli fiber には構造対象を保つ自己同値の作用が入る。
その群oidと軌道集合の正確な区別は第\ref{sec:marked-moduli}節で与える。
\end{theorem}
\begin{proof}
基礎圏を同定すると構造対象をその圏内に輸送できる。
二つの同定の差は exact 自己同値なので、同定を忘れるとその軌道で同一視する。
例外的部分は許容線束の定義そのものである。
\end{proof}
